\begin{lemma}[High-order decay from diagonal vanishing and a H\"older top derivative]
Let $u\colon\mathrm{Plane}\to\mathbb{R}$, let $x_0\in\mathrm{Plane}$, let $\rho>0$, let
$k\in\mathbb{N}$ with $k\ge1$, let $\beta>0$ and let $C\ge0$.  Assume:
\begin{enumerate}
\item $u$ is of class $C^{k}$ on $B_\rho(x_0)$;
\item $D^{j}u(x_0)[w,\ldots,w]=0$ for every $j$ with $0\le j\le k$ and every
$w\in\mathrm{Plane}$, where $D^{j}u(x_0)$ is the $j$-th Fr\'echet derivative of $u$ at $x_0$
evaluated on $j$ copies of $w$;
\item $\bigl\lVert D^{k}u(x)-D^{k}u(y)\bigr\rVert\le C\,\lvert x-y\rvert^{\beta}$ for all
$x,y\in B_\rho(x_0)$, where $\lVert\cdot\rVert$ is the operator norm on continuous $k$-linear
forms on $\mathrm{Plane}$.
\end{enumerate}
Then
\[
  \lvert u(x)\rvert\le\frac{C}{k!}\,\lvert x-x_0\rvert^{\,k+\beta}
  \qquad\text{for every }x\in B_\rho(x_0),
\]
the exponent $k+\beta$ being a real power of the nonnegative real number $\lvert x-x_0\rvert$.
\end{lemma}
\begin{proof}
Fix $x\in B_\rho(x_0)$ and put
\[
  e:=x-x_0,\qquad \gamma(s):=x_0+s\,e,\qquad h(s):=u(\gamma(s))\quad(s\in\mathbb{R}),
\]
so $\gamma(0)=x_0$, $\gamma(1)=x$ and $h(1)=u(x)$.  Since $x\in B_\rho(x_0)$ we have
$\lvert e\rvert<\rho$.  Note $k+\beta>0$, because $k\ge1$ and $\beta>0$.

\medskip\noindent\emph{Case 1: $x=x_0$.}
Then $e=0$.  Hypothesis (2) with $j=0$ and $w=0$ reads $D^{0}u(x_0)[\;]=0$; the $0$-th
Fr\'echet derivative of $u$ at $x_0$ is, by clause (1) of the background proposition on the
recursion for the iterated Fr\'echet derivative, the constant $0$-linear form with value
$u(x_0)$, so this says $u(x_0)=0$.  Hence the left-hand side is $\lvert u(x)\rvert=0$.  The right-hand side
is $\tfrac{C}{k!}\lvert 0\rvert^{\,k+\beta}=\tfrac{C}{k!}\cdot0^{\,k+\beta}=0$, because
$k+\beta>0$ and $0^{a}=0$ for $a>0$.  So the inequality holds, both sides being $0$.

\medskip\noindent\emph{Case 2: $x\neq x_0$, so $0<\lvert e\rvert<\rho$.}
Set
\[
  I:=\{s\in\mathbb{R}:\gamma(s)\in B_\rho(x_0)\}
   =\{s\in\mathbb{R}:\lvert s\rvert\,\lvert e\rvert<\rho\}
   =\left(-\frac{\rho}{\lvert e\rvert},\;\frac{\rho}{\lvert e\rvert}\right),
\]
using $\lvert\gamma(s)-x_0\rvert=\lvert s\,e\rvert=\lvert s\rvert\lvert e\rvert$.  Since
$0<\lvert e\rvert<\rho$ we have $\rho/\lvert e\rvert>1$, so $I$ is an open interval containing
the closed interval $[0,1]$.

\emph{Step 1: $h$ is $C^{k}$ on $I$.}  The ball $B_\rho(x_0)$ is open and $u$ is of class
$C^{k}$ on it by hypothesis (1); the path $\gamma$ is affine, hence $C^\infty$, with
$\gamma'(s)=e$.  By the chain-rule background proposition, $h=u\circ\gamma$ is of class $C^{k}$
on the open set $I=\gamma^{-1}(B_\rho(x_0))$.

\emph{Step 2: the one-variable derivatives of $h$ at $0$ all vanish up to order $k$.}
Let $0\le j\le k$.  Since $u$ is of class $C^{k}$ on $B_\rho(x_0)$ it is of class $C^{j}$
there, and $\gamma(0)=x_0\in B_\rho(x_0)$.  So \noderef{IteratedDerivAlongRay}, applied with
the open set $U=B_\rho(x_0)$, the function $u$, the order $j$, the vectors $x_0$ and $e$ and
the parameter $t=0$, gives
\begin{equation}\label{eq:tvd-at0}
  h^{(j)}(0)=D^{j}u(x_0)[e,\ldots,e]\qquad(0\le j\le k).
\end{equation}
By hypothesis (2) with $w=e$, the right-hand side of \eqref{eq:tvd-at0} is $0$.  Hence
\begin{equation}\label{eq:tvd-vanish}
  h^{(j)}(0)=0\qquad\text{for every }0\le j\le k .
\end{equation}

\emph{Step 3: Taylor's theorem of order $k-1$ along the ray.}
Apply the Taylor--Lagrange background proposition with $n:=k-1$ (legitimate since $k\ge1$, so
$n\in\mathbb{N}$ and $n+1=k$), with $a:=0$, $b:=1$ (these are distinct) and with the function
$h$, which by Step 1 is of class $C^{n+1}=C^{k}$ on the open interval $I\supseteq[0,1]$.  It
produces $\xi\in(0,1)$ with
\[
  h(1)=\sum_{j=0}^{k-1}\frac{h^{(j)}(0)}{j!}\,(1-0)^{j}
       +\frac{h^{(k)}(\xi)}{k!}\,(1-0)^{k}
      =\sum_{j=0}^{k-1}\frac{h^{(j)}(0)}{j!}+\frac{h^{(k)}(\xi)}{k!} .
\]
By \eqref{eq:tvd-vanish} every term of the sum vanishes, so
\begin{equation}\label{eq:tvd-rem}
  u(x)=h(1)=\frac{h^{(k)}(\xi)}{k!} .
\end{equation}

\emph{Step 4: estimating the remainder.}
Since $\xi\in(0,1)\subseteq I$, the point $\gamma(\xi)=x_0+\xi\,e$ lies in $B_\rho(x_0)$.
Applying \noderef{IteratedDerivAlongRay} once more, now with order $k$ and parameter
$t=\xi$ (again $u$ is of class $C^{k}$ on the open set $B_\rho(x_0)$ and
$\gamma(\xi)\in B_\rho(x_0)$), we get
\[
  h^{(k)}(\xi)=D^{k}u(\gamma(\xi))[e,\ldots,e].
\]
Combining this with the case $j=k$ of \eqref{eq:tvd-at0} and \eqref{eq:tvd-vanish}, which give
$D^{k}u(x_0)[e,\ldots,e]=h^{(k)}(0)=0$, and using that evaluation at the fixed family
$(e,\ldots,e)$ is linear in the multilinear form,
\[
  h^{(k)}(\xi)=D^{k}u(\gamma(\xi))[e,\ldots,e]-D^{k}u(x_0)[e,\ldots,e]
             =\bigl(D^{k}u(\gamma(\xi))-D^{k}u(x_0)\bigr)[e,\ldots,e].
\]
The operator-norm background proposition therefore yields
\[
  \bigl\lvert h^{(k)}(\xi)\bigr\rvert
  \le\bigl\lVert D^{k}u(\gamma(\xi))-D^{k}u(x_0)\bigr\rVert\;\lvert e\rvert^{k}.
\]
Both $\gamma(\xi)$ and $x_0$ lie in $B_\rho(x_0)$, so hypothesis (3) applies to the pair
$(\gamma(\xi),x_0)$ and gives
\[
  \bigl\lVert D^{k}u(\gamma(\xi))-D^{k}u(x_0)\bigr\rVert
  \le C\,\lvert\gamma(\xi)-x_0\rvert^{\beta}
  = C\,\bigl(\xi\,\lvert e\rvert\bigr)^{\beta}
  = C\,\xi^{\beta}\lvert e\rvert^{\beta}
  \le C\,\lvert e\rvert^{\beta},
\]
where we used $\lvert\gamma(\xi)-x_0\rvert=\lvert\xi\,e\rvert=\xi\lvert e\rvert$ (valid since
$\xi>0$), the multiplicativity $(ab)^{\beta}=a^{\beta}b^{\beta}$ of real powers of nonnegative
reals, and finally $\xi^{\beta}\le1$, which holds because $0<\xi<1$ and $\beta>0$, together
with $C\ge0$ and $\lvert e\rvert^{\beta}\ge0$.  Hence
\[
  \bigl\lvert h^{(k)}(\xi)\bigr\rvert
  \le C\,\lvert e\rvert^{\beta}\,\lvert e\rvert^{k}
  = C\,\lvert e\rvert^{\,k+\beta},
\]
the last equality being the addition rule for real exponents, $t^{a}t^{b}=t^{a+b}$, applied to
$t=\lvert e\rvert>0$.

\emph{Step 5: conclusion.}  Taking absolute values in \eqref{eq:tvd-rem} and inserting the
bound of Step 4,
\[
  \lvert u(x)\rvert=\frac{\bigl\lvert h^{(k)}(\xi)\bigr\rvert}{k!}
  \le\frac{C}{k!}\,\lvert e\rvert^{\,k+\beta}
  =\frac{C}{k!}\,\lvert x-x_0\rvert^{\,k+\beta},
\]
since $k!>0$.  This is the asserted bound, and the two cases together cover every
$x\in B_\rho(x_0)$.
\end{proof}
